\chapter{Набор инструкций}

Варианты программы для вычисления чисел Фибоначчи

\section{Без дешифратора}

Инструкции:

add c, a, b
mov a, b
mov b, c
nop

Переход не нужен, т.к. ПЗУ зацикливается.

Начальное значение нужно задать вручную

Два бита - тип инструкции

00 - nop
01 - add
10 - mov

Шесть битов - источник и приемник для mov:
001 - a
010 - b
100 - c

Либо просто 3 бита для трех видов инструкций
001 - add c, a, b
010 - mov a, b
100 - mov b, c

% LD_C = INSN_C & Cx
% SEL_C = INSN_C & Cy

Использовать AND или шинные формирователи: сигнал Cx подается на SEL формирователя, а с формирователя DATA идет на LD регистра.

Можно использовать шинный формирователь в обратную сторону, чтобы выбирать из нескольких сигналов.

Регистры
 PC - счётчик инструкций
 T - временный выходной регистр АЛУ
 I - загруженная инструкция

Выборка инструкции
 1 подключить PC
   сбросить I
 2 подключить ROM
   загрузить I
   сбросить T
 3 подключить PC к сумматору
   загрузить T
 4 сбросить PC
 5 скопировать T->PC

Декодирование инструкции
 1 подключить операнды к АЛУ
   сбросить T
 2 записать T
 3 сбросить целевой регистр
 4 загрузить целевой регистр

Если обновлять PC в конце, можно обойтись без I

1 подключить PC
2 сбросить Dst
3 выбрать Src
4 выбрать Src
  выбрать Dst


\section{Программа для вычисления чисел Фибоначчи с переходом}

Фиксированные операнды и аккумулятор

mov a, 0
mov b, 1
L:
add c, a, b
mov a, b
mov b, c
jmp L

\section{АЛУ с выбором регистра}
